% GENERATED FILE. Edit canonical lessons, then run npm run generate:books.
% canonical-source-hash: fnv1a64:f048042f5c4f4732
% canonical-lessons: SA-C517-caturangam, SA-C517-prahelika, SA-C517-kridakah, SA-C517-gulika, SA-C517-putrika, SA-R517-review-510-517

\chapter{Games at Home}
\label{ch:sa-games-at-home}
\hlchaptermodality{517}

\begin{chapteropening}
\textbf{By the end of this chapter:} \emph{I can talk about board games, riddles and guessing.}

Complete the last lesson of chapter 517: I can talk about board games, riddles and guessing.
\end{chapteropening}

\section[\texorpdfstring{\sk{चतुरंगम्} (caturaṅgam)}{चतुरंगम् (caturaṅgam)}]{\sk{चतुरंगम्} (caturaṅgam) — chess}

\label{lesson:SA-C517-caturangam}



\begin{quote}
Before the new one: say the Sanskrit for a collection, then the Sanskrit for a liking.
\end{quote}

\subsection*{You'll want to know: \sk{चतुरंगम्}}
\textbf{\sk{चतुरंगम्}} — \emph{caturaṅgam} — \textquotedblleft{}chess\textquotedblright{}.

\begin{quote}
\sk{अरुणः} \sk{मीरा} \sk{च} \sk{चतुरंगं} \sk{क्रीडतः।}
\end{quote}

\begin{quote}
\emph{aruṇaḥ mīrā ca caturaṅgaṃ krīḍataḥ.}
\end{quote}

\begin{quote}
Arun and Mira are playing chess.
\end{quote}

A neuter noun.

\begin{cousinweb}
\textbf{\sk{चतुरंगम्}} (\emph{caturaṅgam}) is built from \emph{catur}, \textquotedblleft{}four\textquotedblright{}, and \emph{aṅga}, \textquotedblleft{}a limb\textquotedblright{}: the four arms of an old army; MW gives \emph{caturaṅga} also as a kind of chess played by four (MW).
\end{cousinweb}

\begin{grammarlens}[title={What you've built}]
One more word for describing what you lived and felt.
\end{grammarlens}

\subsection*{Your turn}
\begin{itemize}
\raggedright
  \item \emph{Say it:} \emph{caturaṅgam}
  \item \emph{Say it:} \emph{caturaṅgam}, once more
  \item \emph{Say it:} \emph{aruṇaḥ mīrā ca caturaṅgaṃ krīḍataḥ.} — Arun and Mira are playing chess.
  \item \emph{Recall:} say the Sanskrit for a collection, then the Sanskrit for a liking, then say \emph{caturaṅgam} again
\end{itemize}

\begin{tcolorbox}[breakable,colback=teal!4,colframe=teal!35!black,title={Before you move on}]
What does \sk{चतुरंगम्} mean? (\textquotedblleft{}Chess\textquotedblright{}.) Say it once more.
\end{tcolorbox}
\section[\texorpdfstring{\sk{प्रहेलिका} (prahelikā)}{प्रहेलिका (prahelikā)}]{\sk{प्रहेलिका} (prahelikā) — a riddle}

\label{lesson:SA-C517-prahelika}



\begin{quote}
Before the new one: say the Sanskrit for a liking, then the Sanskrit for chess.
\end{quote}

\subsection*{You'll want to know: \sk{प्रहेलिका}}
\textbf{\sk{प्रहेलिका}} — \emph{prahelikā} — \textquotedblleft{}a riddle\textquotedblright{}.

\begin{quote}
\sk{पितामहः} \sk{प्रहेलिकां} \sk{पृच्छति।}
\end{quote}

\begin{quote}
\emph{pitāmahaḥ prahelikāṃ pṛcchati.}
\end{quote}

\begin{quote}
Grandfather asks a riddle.
\end{quote}

A feminine noun.

\begin{cousinweb}
\textbf{\sk{प्रहेलिका}} (\emph{prahelikā}) is the word MW records for a riddle, an enigma; MW gives no derivation, so its origin is uncertain (MW).
\end{cousinweb}

\begin{grammarlens}[title={What you've built}]
One more word for describing what you lived and felt.
\end{grammarlens}

\subsection*{Your turn}
\begin{itemize}
\raggedright
  \item \emph{Say it:} \emph{prahelikā}
  \item \emph{Say it:} \emph{prahelikā}, once more
  \item \emph{Say it:} \emph{pitāmahaḥ prahelikāṃ pṛcchati.} — Grandfather asks a riddle.
  \item \emph{Recall:} say the Sanskrit for a liking, then the Sanskrit for chess, then say \emph{prahelikā} again
\end{itemize}

\begin{tcolorbox}[breakable,colback=teal!4,colframe=teal!35!black,title={Before you move on}]
What does \sk{प्रहेलिका} mean? (\textquotedblleft{}A riddle\textquotedblright{}.) Say it once more.
\end{tcolorbox}
\section[\texorpdfstring{\sk{क्रीडकः} (krīḍakaḥ)}{क्रीडकः (krīḍakaḥ)}]{\sk{क्रीडकः} (krīḍakaḥ) — a player}

\label{lesson:SA-C517-kridakah}



\begin{quote}
Before the new one: say the Sanskrit for chess, then the Sanskrit for a riddle.
\end{quote}

\subsection*{You'll want to know: \sk{क्रीडकः}}
\textbf{\sk{क्रीडकः}} — \emph{krīḍakaḥ} — \textquotedblleft{}a player\textquotedblright{}.

\begin{quote}
\sk{सः} \sk{उत्तमः} \sk{क्रीडकः} \sk{अस्ति।}
\end{quote}

\begin{quote}
\emph{saḥ uttamaḥ krīḍakaḥ asti.}
\end{quote}

\begin{quote}
He is a good player.
\end{quote}

A masculine noun.

\begin{cousinweb}
\textbf{\sk{क्रीडकः}} (\emph{krīḍakaḥ}) comes from the root √krīḍ, \textquotedblleft{}to play\textquotedblright{}, which also gives \emph{krīḍati}, with the agent ending \emph{-aka}; MW gives \emph{krīḍaka} as one who sports, a player (MW).
\end{cousinweb}

\begin{grammarlens}[title={What you've built}]
One more word for describing what you lived and felt.
\end{grammarlens}

\subsection*{Your turn}
\begin{itemize}
\raggedright
  \item \emph{Say it:} \emph{krīḍakaḥ}
  \item \emph{Say it:} \emph{krīḍakaḥ}, once more
  \item \emph{Say it:} \emph{saḥ uttamaḥ krīḍakaḥ asti.} — He is a good player.
  \item \emph{Recall:} say the Sanskrit for chess, then the Sanskrit for a riddle, then say \emph{krīḍakaḥ} again
\end{itemize}

\begin{tcolorbox}[breakable,colback=teal!4,colframe=teal!35!black,title={Before you move on}]
What does \sk{क्रीडकः} mean? (\textquotedblleft{}A player\textquotedblright{}.) Say it once more.
\end{tcolorbox}
\section[\texorpdfstring{\sk{गुलिका} (gulikā)}{गुलिका (gulikā)}]{\sk{गुलिका} (gulikā) — a marble; a small ball}

\label{lesson:SA-C517-gulika}



\begin{quote}
Before the new one: say the Sanskrit for a riddle, then the Sanskrit for a player.
\end{quote}

\subsection*{You'll want to know: \sk{गुलिका}}
\textbf{\sk{गुलिका}} — \emph{gulikā} — \textquotedblleft{}a marble; a small ball\textquotedblright{}.

\begin{quote}
\sk{बालः} \sk{उद्याने} \sk{गुलिकया} \sk{क्रीडति।}
\end{quote}

\begin{quote}
\emph{bālaḥ udyāne gulikayā krīḍati.}
\end{quote}

\begin{quote}
The boy plays marbles in the park.
\end{quote}

A feminine noun.

\begin{cousinweb}
\textbf{\sk{गुलिका}} (\emph{gulikā}) is the word MW records, the same as \emph{guḍikā}, for a ball, a small globule, and a ball to play with; MW gives no derivation, so its origin is uncertain (MW).
\end{cousinweb}

\begin{grammarlens}[title={What you've built}]
One more word for describing what you lived and felt.
\end{grammarlens}

\subsection*{Your turn}
\begin{itemize}
\raggedright
  \item \emph{Say it:} \emph{gulikā}
  \item \emph{Say it:} \emph{gulikā}, once more
  \item \emph{Say it:} \emph{bālaḥ udyāne gulikayā krīḍati.} — The boy plays marbles in the park.
  \item \emph{Recall:} say the Sanskrit for a riddle, then the Sanskrit for a player, then say \emph{gulikā} again
\end{itemize}

\begin{tcolorbox}[breakable,colback=teal!4,colframe=teal!35!black,title={Before you move on}]
What does \sk{गुलिका} mean? (\textquotedblleft{}A marble; a small ball\textquotedblright{}.) Say it once more.
\end{tcolorbox}
\section[\texorpdfstring{\sk{पुत्रिका} (putrikā)}{पुत्रिका (putrikā)}]{\sk{पुत्रिका} (putrikā) — a doll}

\label{lesson:SA-C517-putrika}



\begin{quote}
Before the new one: say the Sanskrit for a player, then the Sanskrit for a marble.
\end{quote}

\subsection*{You'll want to know: \sk{पुत्रिका}}
\textbf{\sk{पुत्रिका}} — \emph{putrikā} — \textquotedblleft{}a doll\textquotedblright{}.

\begin{quote}
\sk{बालिका} \sk{पुत्रिकया} \sk{क्रीडति।}
\end{quote}

\begin{quote}
\emph{bālikā putrikayā krīḍati.}
\end{quote}

\begin{quote}
The girl plays with a doll.
\end{quote}

A feminine noun.

\begin{cousinweb}
\textbf{\sk{पुत्रिका}} (\emph{putrikā}) is built on \emph{putra}, \textquotedblleft{}a son, a child\textquotedblright{}, with the ending \emph{-ikā}; MW gives \emph{putrikā} as a daughter, and as a puppet, a doll (MW).
\end{cousinweb}

\begin{grammarlens}[title={What you've built}]
One more word for describing what you lived and felt.
\end{grammarlens}

\subsection*{Your turn}
\begin{itemize}
\raggedright
  \item \emph{Say it:} \emph{putrikā}
  \item \emph{Say it:} \emph{putrikā}, once more
  \item \emph{Say it:} \emph{bālikā putrikayā krīḍati.} — The girl plays with a doll.
  \item \emph{Recall:} say the Sanskrit for a player, then the Sanskrit for a marble, then say \emph{putrikā} again
\end{itemize}

\begin{tcolorbox}[breakable,colback=teal!4,colframe=teal!35!black,title={Before you move on}]
What does \sk{पुत्रिका} mean? (\textquotedblleft{}A doll\textquotedblright{}.) Say it once more.
\end{tcolorbox}
\section[(dialogue)]{Looking back: from Clearing the Clutter to Games at Home}

\label{lesson:SA-R517-review-510-517}



\begin{quote}
Say the Sanskrit for rubbish and clears away.
\end{quote}

\subsection*{Your turn}
Say these aloud:

\begin{itemize}
\raggedright
  \item rubbish, clears away, puts in order, a pile, a thorough cleaning
  \item firewood, makes hot, takes down, cooking, roasts
  \item a knife, cuts up, grinds to powder, a winnowing basket, combines
  \item churns, a churning stick, ghee, strains, oily
  \item serves food, feeds, dines, leftovers, a lid
  \item a flavour, tasty, raw, ripe, sucks
  \item entertainment, fun, amuses, a collection, a liking
  \item chess, a riddle, a player, a marble, a doll
\end{itemize}

\begin{tcolorbox}[breakable,colback=teal!4,colframe=teal!35!black,title={Before you move on}]
Rubbish? (\textbf{\sk{अवकरः}}, \emph{avakaraḥ}.) Clears away? (\textbf{\sk{उत्सारयति}}, \emph{utsārayati}.) Puts in order? (\textbf{\sk{व्यवस्थापयति}}, \emph{vyavasthāpayati}.) A pile? (\textbf{\sk{पुञ्जः}}, \emph{puñjaḥ}.) A thorough cleaning? (\textbf{\sk{सम्मार्जनम्}}, \emph{sammārjanam}.) Firewood? (\textbf{\sk{काष्ठम्}}, \emph{kāṣṭham}.) Makes hot? (\textbf{\sk{तापयति}}, \emph{tāpayati}.) Takes down? (\textbf{\sk{उत्तारयति}}, \emph{uttārayati}.) Cooking? (\textbf{\sk{पाकः}}, \emph{pākaḥ}.) Roasts? (\textbf{\sk{भृज्जति}}, \emph{bhṛjjati}.) A knife? (\textbf{\sk{छुरिका}}, \emph{churikā}.) Cuts up? (\textbf{\sk{कृन्तति}}, \emph{kṛntati}.) Grinds to powder? (\textbf{\sk{चूर्णयति}}, \emph{cūrṇayati}.) A winnowing basket? (\textbf{\sk{शूर्पम्}}, \emph{śūrpam}.) Combines? (\textbf{\sk{संयोजयति}}, \emph{saṃyojayati}.) Churns? (\textbf{\sk{मथ्नाति}}, \emph{mathnāti}.) A churning stick? (\textbf{\sk{मन्थः}}, \emph{manthaḥ}.) Ghee? (\textbf{\sk{घृतम्}}, \emph{ghṛtam}.) Strains? (\textbf{\sk{गालयति}}, \emph{gālayati}.) Oily? (\textbf{\sk{स्निग्धः}}, \emph{snigdhaḥ}.) Serves food? (\textbf{\sk{परिवेषयति}}, \emph{pariveṣayati}.) Feeds? (\textbf{\sk{भोजयति}}, \emph{bhojayati}.) Dines? (\textbf{\sk{अश्नाति}}, \emph{aśnāti}.) Leftovers? (\textbf{\sk{उच्छिष्टम्}}, \emph{ucchiṣṭam}.) A lid? (\textbf{\sk{पिधानम्}}, \emph{pidhānam}.) A flavour? (\textbf{\sk{स्वादः}}, \emph{svādaḥ}.) Tasty? (\textbf{\sk{स्वादुः}}, \emph{svāduḥ}.) Raw? (\textbf{\sk{आमः}}, \emph{āmaḥ}.) Ripe? (\textbf{\sk{पक्वः}}, \emph{pakvaḥ}.) Sucks? (\textbf{\sk{चूषति}}, \emph{cūṣati}.) Entertainment? (\textbf{\sk{मनोरञ्जनम्}}, \emph{manorañjanam}.) Fun? (\textbf{\sk{हास्यम्}}, \emph{hāsyam}.) Amuses? (\textbf{\sk{विनोदयति}}, \emph{vinodayati}.) A collection? (\textbf{\sk{संग्रहः}}, \emph{saṅgrahaḥ}.) A liking? (\textbf{\sk{रुचिः}}, \emph{ruciḥ}.) Chess? (\textbf{\sk{चतुरंगम्}}, \emph{caturaṅgam}.) A riddle? (\textbf{\sk{प्रहेलिका}}, \emph{prahelikā}.) A player? (\textbf{\sk{क्रीडकः}}, \emph{krīḍakaḥ}.) A marble? (\textbf{\sk{गुलिका}}, \emph{gulikā}.) A doll? (\textbf{\sk{पुत्रिका}}, \emph{putrikā}.)
\end{tcolorbox}
